\section{Quantitative Shioda--Katsura suppliers for all products of
Fermat-twelve
hypersurfaces}\label{sec:sf-fermat12-hypersurfaces}

Author/checker: GPT-6 Astra Pro; inherited effort not exposed;
HODGE\_DIVISION /root/sixfolds. Date: 2026-11-21. New full proof; root whole-proof independent audit PASS on2026-11-21,
including an actual rerun of all423 finite degree cases. Reviewer:
GPT-6 Astra Pro; inherited effort not exposed; HODGE\_DIVISION /root. This extends the independently checked
SF-F1201--SF-F1204.

Write \(F^n=F_{12}^n\subset\mathbf P^{n+1}\), with \(\sum x_i^{12}=0\),
and give every product its Segre hyperplane class, the sum of the
hyperplane classes of its factors. A zero-dimensional Fermat factor is
treated componentwise; it consists of twelve reduced points.

For \(T\ge1\), define the following explicit integers: \[
 \begin{aligned}
 P_T&=12^T T!,\\
 A_T&=2^{2T}P_T^2,\\
 M_T&=2(2T+1)\,3^T A_T^{\,2T+2},\\
 U_T&=440^T T!\,144^{55T},\\
 \mathcal H_{12}(T)&=
 U_T\prod_{j=2}^{T}M_j^{\,j-1}.
 \end{aligned}
 \tag{SKQ.1}
\] Put \(\mathcal H_{12}(0)=1\); an empty product is one.

\textbf{Theorem.} For every product \(X=\prod_\nu F^{d_\nu}\), with
\(d_\nu\ge1\) and total dimension \(T=\sum d_\nu\), and every
codimension, the rational Hodge classes of \(X\) are spanned by a finite
set of effective integral algebraic cycles, each having Segre degree at
most \(\mathcal H_{12}(T)\).

This is a numerical bound for actual effective cycles, not a bound for
formal Hodge tensors alone. The assertion concerns rational spanning. No
saturation of the integral Hodge lattice is claimed.

\subsection{SF-SKQ01. A bounded effective supplier on every product of
curves}\label{SF-SKQ01}

Let \(C=F^1\), with plane hyperplane class \(h\), and let
\(a:C\hookrightarrow J(C)\) be an Abel embedding. Under the Abel
identification on \(H^1\), the canonical theta class is
\(\sum_{j=1}^{55}\alpha_j\beta_j\) for a symplectic integral basis with
\(\int_C\alpha_j\beta_j=1\). Its pullback therefore has degree55, while
the plane hyperplane class has degree12: \[
 a^*\Theta=55[\mathrm{pt}],\qquad h=12[\mathrm{pt}]
 \quad\hbox{in }H^2(C,\mathbf Q).
\] For \(a_T:C^T\to J(C)^T\), let \(H_C=\sum\operatorname{pr}_j^*h\).
Then \[
 a_T^*\Theta_T=\frac{55}{12}H_C,\qquad
 \int_{C^T}H_C^T=T!\,12^T.
 \tag{SKQ.2}
\]

The pullback \(a_T^*:H^*(J(C)^T,\mathbf Q)\to H^*(C^T,\mathbf Q)\) is
surjective as a graded Hodge structure and admits a rational Hodge
section. On a single factor, use the inverse isomorphism in degree one,
the identity in degree zero, and the map
\([\mathrm{pt}]\mapsto\Theta/55\) in degree two. Tensor these sections
over all factors. Consequently every rational Hodge class of \(C^T\) has
a rational Hodge lift to \(J(C)^T\).

Use the effective cycle suppliers \(Z=(\Phi_T)_*Y\) of SF-F1204 in
codimension \(k\le T\). The cycles \(Y\) are simultaneous kernels of
elliptic homomorphisms on \(E_i^{27T}\times E_\omega^{28T}\). They are
finite unions of abelian subvarieties translated by torsion points.
Their finite-isogeny images \(Z\) have the same form, with positive
integral multiplicities.

A general translation of every component of \(Z\) meets \(a_T(C^T)\)
properly or not at all. Here no general moving bound is needed. For an
abelian subvariety \(A\) of codimension \(k\), compose \(a_T\) with the
quotient \(J(C)^T\to J(C)^T/A\). A general fibre either is empty, when
the map is not dominant, or has dimension \(T-k\), by generic
smoothness. The same choice works for the finitely many components. Thus
the pullback of a sufficiently general translate of \(Z\) is an
effective integral cycle \(W\) of codimension \(k\), possibly zero,
representing \(a_T^*[Z]\).

There is also a degree bound substantially smaller than applying a
general intersection estimate in the Jacobian. On the elliptic source of
\(\Phi_T\), each generating divisor satisfies \(D_j\le8L_T\), as
semipositive translation-invariant forms. Expanding one factor at a time
gives \[
 D_{j_1}\cdots D_{j_k}\le8^kL_T^k
\] in the cone of strongly positive \((k,k)\)-forms. Pushforward by the
finite isogeny preserves this inequality, and \[
 (\Phi_T)_*L_T^k=(\deg\Phi_T)\Theta_T^k,\qquad
 \deg\Phi_T\le144^{55T}.
\] Pulling back to the curve product and integrating against the
positive form \(H_C^{T-k}\) therefore yields \[
 \begin{aligned}
 \deg_{H_C}W
 &\le8^k144^{55T}
       \int_{C^T}(a_T^*\Theta_T)^kH_C^{T-k}\\
 &=8^k55^k12^{T-k}T!\,144^{55T}\\
 &\le440^T T!\,144^{55T}=U_T.
 \end{aligned}
 \tag{SKQ.3}
\] Translations do not alter these cohomology classes. Since the
original Jacobian classes span and \(a_T^*\) has a Hodge section, the
resulting effective curve-product cycles span every rational Hodge
class.

The degree-zero case uses the whole curve product; the top-codimension
case is allowed. Zero pullbacks are discarded. This proves the base of
the induction in all codimensions simultaneously.

\subsection{SF-SKQ02. The two explicit correspondences in one induction
step}\label{SF-SKQ02}

Suppose \[
 X=F^n\times R,\qquad n\ge2,\qquad \dim X=T,
\] where \(R\) is a product of positive-dimensional Fermat
hypersurfaces. Put \[
 Z=F^{n-1}\times C\times R,\qquad
 B=F^{n-2}\times F^0\times R.
\] Here \(B\subset Z\) is cut out by the last homogeneous coordinates
\(x_n=0\), \(y_2=0\), so it is smooth of codimension two, possibly
disconnected.

Fix \(\epsilon^{12}=-1\). The rational map on the first two factors is
\[
 ((x_0:\cdots:x_n),(y_0:y_1:y_2))
 \longmapsto
 (x_0y_2:\cdots:x_{n-1}y_2:
       \epsilon x_ny_0:\epsilon x_ny_1).
 \tag{SKQ.4}
\] Its coordinates have bidegree \((1,1)\), its base ideal is the ideal
of \(B\), and its generic degree is12. For the last assertion, on a
target point where the sum of the first \(n\) twelfth powers is nonzero,
normalizing \(y_2=1\) reconstructs \(x_n\) from a nonzero prescribed
twelfth power, giving twelve distinct preimages.

Let \(\pi:\widetilde Z=\operatorname{Bl}_B Z\to Z\). Locally the base
ideal is generated by the two regular parameters \(x_n,y_2\); the blowup
therefore resolves (SKQ.4), giving a morphism \(f:\widetilde Z\to X\),
of generic degree12, and the identity on \(R\). This is the concrete
Shioda--Katsura construction: Lemmas1.1--1.2, pp98--99, and Theorem1.7,
pp101--102, of their published1979 paper.

Write \(E\) for the exceptional divisor,
\(j:E\hookrightarrow\widetilde Z\), and \(p:E\to B\). For a
codimension-two smooth blowup, the rational Hodge decomposition is \[
 H^m(\widetilde Z,\mathbf Q)
 =
 \pi^*H^m(Z,\mathbf Q)
 \oplus j_*p^*H^{m-2}(B,\mathbf Q)(-1).
 \tag{SKQ.5}
\] The maps displayed here are the actual algebraic pullback and Gysin
maps. The inverse components are \(\pi_*\) and \(-p_*j^*\): one uses
\(j^*j_*\eta=c_1(\mathcal O_E(-1))\eta\) and the projective-line bundle
formula. The blowup formula follows by excision and the cohomology of
the exceptional projective bundle; the same geometric formula is the
input of Shioda--Katsura Lemma2.1 and Proposition2.4, pp103--104. Over
\(\mathbf C\) these algebraic maps respect the stated Hodge structures.

For every rational Hodge class \(\alpha\) on \(X\), \[
 12\alpha=f_*f^*\alpha.
\] Decompose \(f^*\alpha\) by (SKQ.5). Thus the entire Hodge space of
\(X\) is spanned by the images of Hodge classes on \(Z\) and on the
connected components \(B_\lambda\) under the correspondences \[
 \Gamma_0=(\pi,f)_*[\widetilde Z]\subset Z\times X,\qquad
 \Gamma_\lambda=(p,fj)_*[E|_{B_\lambda}]
                  \subset B_\lambda\times X.
 \tag{SKQ.6}
\] The first acts by \(f_*\pi^*\); the second acts by \(f_*j_*p^*\),
increasing codimension by one. This includes every exceptional and Tate
contribution. No claim about a primitive summand is substituted for the
full decomposition.

\subsection{SF-SKQ03. Degrees of the main and exceptional
correspondences}\label{SF-SKQ03}

Every product of positive-dimensional degree12 hypersurfaces of total
dimension \(T\) has Segre degree at most \[
 P_T=12^T T!.
 \tag{SKQ.7}
\] Indeed its exact degree is \(T!\prod_\nu12/d_\nu!\), and there are at
most \(T\) factors. The same bound applies to \(X\) and \(Z\).

Let \(H_Z,H_X\) be their Segre hyperplane classes. The resolved linear
system gives \[
 f^*H_X=\pi^*H_Z-E.
\] Hence the hyperplane class from the product embedding \(Z\times X\),
pulled back to \(\widetilde Z\), is \[
 G=2\pi^*H_Z-E.
\] Both \(G\) and \(2\pi^*H_Z\) are nef. Their difference is the
effective divisor \(E\). Expanding their difference of \(T\)-th powers
shows \[
 \int_{\widetilde Z}G^T\le
 2^T\int_{\widetilde Z}(\pi^*H_Z)^T.
\] Each term in the difference is an effective divisor intersected with
nef classes, hence has nonnegative degree. Since \((\pi,f)\) has degree
one onto its graph, \[
 \deg\Gamma_0\le2^T\deg Z\le2^TP_T.
 \tag{SKQ.8}
\]

For completeness the exceptional degree is controlled by its actual
projective bundle, without choosing a resolution of an unspecified
graph. On a connected component \(B_\lambda\), the \(F^0\) coordinate
\(y\) is a fixed point. If \(n=2\), the other \(F^0\) coordinate is
fixed as well. Write \(h\) for the remaining \(F^{n-2}\) hyperplane
class, or zero when that factor is a point, and \(H_R\) for the
hyperplane class of \(R\). Then \[
 H_{B_\lambda}=h+H_R.
\] The exceptional fibres map to the lines joining the point represented
by \(x\) and the fixed point represented by \(y\). Thus their incidence
variety is explicitly \[
 E_\lambda=
 \mathbf P_{\rm lines}\bigl(\mathcal O_{B_\lambda}(-h)
                              \oplus\mathcal O_{B_\lambda}\bigr).
\] Let \(\xi=c_1(\mathcal O_{E_\lambda}(1))\). The join map has
hyperplane class \(\xi\); in particular \(\xi\) is nef. Including the
identity map on \(R\), the pullback of the product hyperplane class from
\(B_\lambda\times X\) is \[
 \xi+h+2H_R\le\xi+2H_{B_\lambda}.
\] The projective bundle relation and pushforward formulas are \[
 \xi(\xi-h)=0,\qquad
 p_*\xi^j=h^{j-1}\quad(j\ge1),\qquad p_*1=0.
\] Therefore \[
 \begin{aligned}
 \deg\Gamma_\lambda
 &\le\int_{E_\lambda}(\xi+2H_{B_\lambda})^{T-1}\\
 &=\sum_{j=1}^{T-1}
       \binom{T-1}{j}2^{T-1-j}
       \int_{B_\lambda}h^{j-1}H_{B_\lambda}^{T-1-j}\\
 &\le(3^{T-1}-2^{T-1})\deg B_\lambda\\
 &\le3^TP_T.
 \end{aligned}
 \tag{SKQ.9}
\] The third line uses \(0\le h\le H_{B_\lambda}\), with both classes
nef. The component \(B_\lambda\) is a product of total dimension
\(T-2\), with any zero-dimensional factors replaced by points, so its
degree is at most \(P_{T-2}\le P_T\). The correspondence map has degree
one onto its incidence image; even the weaker inequality for its
pushforward would suffice.

Thus every correspondence in (SKQ.6) has degree at most \(3^TP_T\) in
the relevant product embedding.

\subsection{SF-SKQ04. The quantitative correspondence-action
bound}\label{SF-SKQ04}

The published intersection estimate used here is Truong, \emph{Relative
dynamical degrees of correspondences over a field of arbitrary
characteristic}, Crelle758(2020),139--182, Lemma2.2 in the author
version, p13. Its full proof uses Roberts' moving construction and
bounds an intersection of effective cycles by a difference of two
effective cycles, each of degree at most \[
 C(V)\,\deg A\,\deg B.
\] The proof gives \(C(V)=m\deg(V)^m\), for \(m=\dim V\). To avoid any
convention about collecting positive and negative terms, use throughout
the larger integer \[
 C_{\rm safe}(V)=2(m+1)\deg(V)^{m+1}.
 \tag{SKQ.10}
\] The estimate extends from irreducible to effective cycles by
bilinearity and summing positive multiplicities. Rational equivalence is
enough here, and therefore the resulting equality also holds in rational
cohomology.

Let \(Y\) be either \(Z\) or a component \(B_\lambda\), and let
\(\Gamma\) be the corresponding cycle in \(Y\times X\). The ambient
product \(V=Y\times X\) is smooth and projective, has dimension at most
\(2T\), and has degree at most \[
 \deg V
 =\binom{\dim Y+T}{\dim Y}\deg Y\deg X
 \le2^{2T}P_T^2=A_T.
\] For an effective input cycle \(W\subset Y\), \[
 \deg(W\times X)\le2^{2T}P_T\deg W.
\] Apply (SKQ.10) to \(\Gamma\cdot(W\times X)\). Since
\(\deg\Gamma\le3^TP_T\), it can be represented as \(Q_+-Q_-\), with both
effective and \[
 \begin{aligned}
 \deg Q_\pm
 &\le
 2(2T+1)A_T^{2T+1}
      (3^TP_T)(2^{2T}P_T\deg W)\\
 &=2(2T+1)3^TA_T^{2T+2}\deg W\\
 &=M_T\deg W.
 \end{aligned}
 \tag{SKQ.11}
\] Push these cycles to \(X\). A component whose image has smaller
dimension pushes to zero. On every other component the projection
formula and the nef inequality \(p_X^*H_X\le H_V\) show that degree does
not increase. Consequently \[
 \Gamma_*[W]=[W_+]-[W_-],\qquad
 W_\pm\ \text{effective integral},\qquad
 \deg W_\pm\le M_T\deg W.
 \tag{SKQ.12}
\] In particular, if bounded-degree effective input classes span the
source Hodge space, then effective outputs of degree at most \(M_T\)
times that bound span the correspondence image. It is harmless that both
signs are retained as separate generators.

No proper-intersection assumption is imposed on the original input and
correspondence. That is exactly what the cited moving lemma supplies,
with the displayed numerical cost.

\subsection{SF-SKQ05. Induction, an explicit recurrence, and the final
bound}\label{SF-SKQ05}

For a product of total dimension \(T\), define its complexity \[
 \mu=\sum_\nu(d_\nu-1)=T-\#\{\hbox{positive-dimensional factors}\}.
\] Then \(0\le\mu\le T-1\). Complexity zero means \(C^T\), already
handled in SF-SKQ01.

Replacing one factor \(F^n\), \(n\ge2\), by \(F^{n-1}\times C\)
preserves \(T\) and decreases \(\mu\) by one. Every exceptional
component has total dimension \(T-2\). Thus (SKQ.6) and (SKQ.12) give
the valid simultaneous recurrence \[
 B(T,0)=U_T,\qquad
 B(T,\mu)\le
 M_T\max\{B(T,\mu-1),\,\mathcal H_{12}(T-2)\},
 \tag{SKQ.13}
\] with the point value \(\mathcal H_{12}(0)=1\). All codimensions are
included at once. Terms whose codimension lies outside the source range
vanish and can simply be omitted.

Iterating at most \(T-1\) times yields \[
 B(T,\mu)\le
 M_T^{T-1}\max\{U_T,\mathcal H_{12}(T-2)\}.
 \tag{SKQ.14}
\] The explicit quantity in (SKQ.1) dominates this expression. Indeed
\(U_T\) increases with \(T\), every \(M_j\ge1\), and \[
 \frac{\mathcal H_{12}(T)}
      {M_T^{T-1}\mathcal H_{12}(T-2)}
 =
 \frac{U_T}{U_{T-2}}\,M_{T-1}^{T-2}\ge1
 \quad(T\ge3);
\] the case \(T=2\) follows directly. This proves the theorem by
induction on total dimension and then complexity.

Only finitely many effective generators are introduced at each stage:
the previous finite generating lists are passed through two
correspondence types, with finitely many exceptional components, and
each intersection is replaced by two effective cycles. The construction
never invokes an uncontrolled resolution, an unbounded regularity
estimate, or an unspecified isogeny degree. Every factor in (SKQ.1) has
been accounted for.

For a nonzero pairing functional on a Hodge space, at least one
generator pairs nontrivially. One may then choose one irreducible
component with nonzero pairing, still of degree at most
\(\mathcal H_{12}(T)\). This is the intended use as a numerical supplier
in the subsequent cyclic-cover and Brauer-correspondence transfer.

\subsection{Published-source and verification
record}\label{published-source-and-verification-record}

Shioda--Katsura is published in Tohoku Mathematical
Journal31(1979),97--115, DOI10.2748/tmj/1178229881. The official J-STAGE
PDF was opened; the rational map and blowup construction on pp98--102
and cohomological formulas on pp103--104 were inspected. The degree
bounds, projective-bundle calculation and full-product recurrence above
are new.

Truong's paper is published in Crelle758(2020),139--182,
DOI10.1515/crelle-2017-0052. Publication was checked at the publisher.
The complete relevant quantitative proof was read in the author copy
arXiv:1605.05049, Lemma2.2, p13. No unverified printed lemma-page
locator is asserted. The larger safe constant (SKQ.10) is used, so the
recurrence does not depend on a sharp collection-of-signs convention.
